\chapter{绪论}
\section{研究背景与问题来源}

等离子体增强化学气相沉积（plasma-enhanced chemical vapor deposition，PECVD）用于半导体薄膜沉积。先进制程中的三维器件和多层薄膜集成提高了沉积均匀性要求，设备同时需要满足高吞吐量需求\citep{appliedprecisioncvd}。组合设备（cluster tool，CT）通过机械手和装载锁连接多个工艺模块，其运行时序受工艺、搬运、清洁和驻留约束共同影响\citep{pan2018survey}。

多空间模块与多指机械手提高了单设备的并行处理和搬运能力\citep{gu2024multifinger}。本文设备配置四指真空机械手（VTR）和双指大气机械手（ATR）：VTR的洁净与脏手各有两个承片位置，ATR可承载两片，旋转腔每次可访问两个工位。因此，以两片晶圆组成的PW对安排腔室加工与成对搬运，同时保留对内两片晶圆各自的供片、对准和回片时刻。

本文考虑采用可复用辅助晶圆保护工位底座的满片加工场景。产品不足时由辅助片补位，清洁时由辅助片占满全部工位。辅助片不计入产出，但借用、搬运和归还占用设备资源。产品从大气侧装载端口经对准器和装载锁进入腔室；单源配置的辅助片也由大气侧供给，双源配置则由真空存储区供给，形成产品与辅助晶圆的混流。

四工位旋转腔支持$4\times1$并行加工和$2\times2$链式加工。前者使四片晶圆同步完成一次处理；后者使产品连续参加两次处理，并以辅助片补足链首和链尾。合成算例中两模式的单片累计处理时长均为180~s，旋转与换片安排使资源占用不同。批次或设备参数变化后需要重新排程，因而研究重点是有限计算时间内的可行解质量。

\begin{figure}[!htbp]
\centering\includegraphics[width=0.72\linewidth]{fig_system_pjob_cn}
\caption{双源混流旋转腔组合设备拓扑与晶圆加工路径}\label{fig:system-zh-new}
\end{figure}

本文采用标称工艺参数构造算例，以分析晶圆组织与资源时序对排程的影响。调度对象为给定任务、设备参数和资源容量下的有限批次离线排程。
\section{研究内容与评价目标}
围绕有限计算时间内的排程质量，本文已开展的研究包括以下三个相互衔接的方面：

\begin{enumerate}
\item 建立以PW对为最小单元的两阶段MILP模型。按任务和来源LP统计产品数量，联合确定晶圆配对、腔室分配、加工周期、清洁和辅助片借还。第一阶段最小化批次完工时间，第二阶段保持该完工上界和主要离散决策，优化等待与节拍。
\item 设计结构化初始解策略。先确定晶圆配对、腔室分配和加工周期，再逐层放开资源使用次序进行优化，为主问题提供较好的初始可行解。
\item 设计变量角色特征提取与结构感知候选割重排。根据候选割中不同类型变量的系数占比进行学习排序，再结合变量类型覆盖程度和候选割相似度选择割子集，在选割数量有限时减少相似割的重复入选。
\end{enumerate}


评价分别考察模型能否恢复满足设备约束的逐片排程、结构化初始解能否降低限时可行上界，以及共享初始解后学习选割能否带来额外改善。清洁规则、加工模式与辅助库存的对照用于解释排程的变化，阶段内对照用于分析等待与节拍。实验同时报告可行上界、有效下界和停止状态。
\section{论文结构}
第一章说明研究背景、研究内容和评价目标。第二章梳理组合设备调度、时序约束与学习选割研究。第三章建立PW对联合模型并定义完工时间优先的两阶段目标。第四章介绍结构化初始解、变量角色特征与候选割重排，以及排程精修和解码流程，并为产品完工瓶颈驱动的原始解改进保留独立研究位置。第五章分析已完成的模型、初始化、选割、清洁与库存实验，展示2片最优排程和16片可行排程，并区分历史100片实例与当前模型的扩展任务。第六章归纳已有研究结论与后续方向。
